\documentclass[10pt,a4paper]{amsart}
\usepackage{amsmath,amssymb,amsthm,graphicx}
% Retained preamble for the Stacks Project chapter snapshot in this directory.
% The chapter's own preamble.tex is a copy of the shared Stacks Project preamble; the block
% below is sliced MECHANICALLY, line for line, from it. Nothing was added, removed or
% redefined; the wrapper class, the xr-hyper/hyperref wiring and the externaldocument list
% are omitted because the assembler supplies the document class.
%
% lines taken from preamble.tex: 3, 5, 9, 11, 12, 156, 157, 158, ... (51 lines total)
%
% package coverage restored from the chapter's own preamble (2026-10-13): xy, xy-2cell

\usepackage[all]{xy}
\xyoption{2cell}
\UseAllTwocells

\newenvironment{reference}{\comment}{\endcomment}
\newenvironment{slogan}{\comment}{\endcomment}
\newenvironment{history}{\comment}{\endcomment}
\theoremstyle{plain}
\newtheorem{theorem}[subsection]{Theorem}
\newtheorem{proposition}[subsection]{Proposition}
\newtheorem{lemma}[subsection]{Lemma}
\theoremstyle{definition}
\newtheorem{definition}[subsection]{Definition}
\newtheorem{example}[subsection]{Example}
\newtheorem{exercise}[subsection]{Exercise}
\newtheorem{situation}[subsection]{Situation}
\theoremstyle{remark}
\newtheorem{remark}[subsection]{Remark}
\newtheorem{remarks}[subsection]{Remarks}
\numberwithin{equation}{subsection}
\def\lim{\mathop{\mathrm{lim}}\nolimits}
\def\colim{\mathop{\mathrm{colim}}\nolimits}
\def\Spec{\mathop{\mathrm{Spec}}}
\def\Hom{\mathop{\mathrm{Hom}}\nolimits}
\def\Ext{\mathop{\mathrm{Ext}}\nolimits}
\def\SheafHom{\mathop{\mathcal{H}\!\mathit{om}}\nolimits}
\def\SheafExt{\mathop{\mathcal{E}\!\mathit{xt}}\nolimits}
\def\Sch{\mathit{Sch}}
\def\Mor{\mathop{\mathrm{Mor}}\nolimits}
\def\Ob{\mathop{\mathrm{Ob}}\nolimits}
\def\Sh{\mathop{\mathit{Sh}}\nolimits}
\def\NL{\mathop{N\!L}\nolimits}
\def\CH{\mathop{\mathrm{CH}}\nolimits}
\def\proetale{{pro\text{-}\acute{e}tale}}
\def\etale{{\acute{e}tale}}
\def\QCoh{\mathit{QCoh}}
\def\Ker{\mathop{\mathrm{Ker}}}
\def\Im{\mathop{\mathrm{Im}}}
\def\Coker{\mathop{\mathrm{Coker}}}
\def\Coim{\mathop{\mathrm{Coim}}}
\DeclareMathSymbol{\boxtimes}{\mathbin}{AMSa}{"02}
\def\QCohstack{\mathcal{QC}\!\mathit{oh}}
\def\Cohstack{\mathcal{C}\!\mathit{oh}}
\def\Spacesstack{\mathcal{S}\!\mathit{paces}}
\def\Quotfunctor{\mathrm{Quot}}
\def\Hilbfunctor{\mathrm{Hilb}}
\def\Curvesstack{\mathcal{C}\!\mathit{urves}}
\def\Polarizedstack{\mathcal{P}\!\mathit{olarized}}
\def\Complexesstack{\mathcal{C}\!\mathit{omplexes}}
\def\Pic{\mathop{\mathrm{Pic}}\nolimits}
\def\Picardstack{\mathcal{P}\!\mathit{ic}}
\def\Picardfunctor{\mathrm{Pic}}
\def\Deformationcategory{\mathcal{D}\!\mathit{ef}}
\usepackage{cleveref}
\crefname{originallemma}{original lemma}{original lemmas}
\Crefname{originallemma}{Original lemma}{Original lemmas}
\makeatletter
\newlabel{equation-completion-functor}{{0880}{source}{}{sourcepointer-1}{}}
\newlabel{equation-completion-functor@cref}{{[originallemma][0][]0880}{[1][0][]source}{}{}{}}
\makeatother
\crefname{originallemma}{original lemma}{original lemmas}
\Crefname{originallemma}{Original lemma}{Original lemmas}
\makeatletter
\newlabel{homology-section-extensions}{{010I}{source}{}{sourcepointer-2}{}}
\newlabel{homology-section-extensions@cref}{{[originallemma][0][]010I}{[1][0][]source}{}{}{}}
\makeatother
\crefname{originallemma}{original lemma}{original lemmas}
\Crefname{originallemma}{Original lemma}{Original lemmas}
\makeatletter
\newlabel{lemma-completion-internal-hom}{{0882}{source}{}{sourcepointer-3}{}}
\newlabel{lemma-completion-internal-hom@cref}{{[originallemma][0][]0882}{[1][0][]source}{}{}{}}
\makeatother
\crefname{originallemma}{original lemma}{original lemmas}
\Crefname{originallemma}{Original lemma}{Original lemmas}
\makeatletter
\newlabel{lemma-existence-easy}{{0889}{source}{}{sourcepointer-4}{}}
\newlabel{lemma-existence-easy@cref}{{[originallemma][0][]0889}{[1][0][]source}{}{}{}}
\makeatother
\title{Stacks, Cohomology of Schemes (coherent.tex): lemma-torsion-hom-ext}
\author{The Stacks Project Authors}
\date{stacks-project:coherent.tex:local-snapshot-20261012; GNU FDL 1.2 or later}
\begin{document}
\maketitle
\noindent\textit{Editorial note.} Statement span is lines 6765-6784 of sources/stacks-coherent/pinned-original-source.tex; proof span is lines 6786-6826, the complete author proof (39 body lines).\par\medskip
\noindent\textit{Editorial note.} Navigation stubs: 4 cited-but-not-included labels are supplied as non-load-bearing pointers with their pinned tags where the snapshot tag index has them.\par\medskip
\noindent\textit{Editorial note.} extra\_packages is empty on purpose: the retained preamble already loads xy with options and a second declaration clashes in the assembler wrapper.\par\medskip
\noindent\textit{Editorial note.} Text is retained verbatim, including the author's own typos and placeholders.\par\medskip
\noindent\textit{Editorial note.} Queue index 47 of queue\_coherent.json (pb=1715, stmt 6765-6784, proof 6786-6826).\par\medskip
\noindent\textit{Editorial note.} co02 wave: conflict-free remaining accept (old-session co01-co08 build did not cover this label; fresh adjudication 2026-10-14 accepted).\par\medskip
\noindent\textit{Editorial note.} gate\_ok: precard\_gate RESULT=OK (tag-based, 2026-10-14); builder raw-interval guard false-positives on problem 376 whose public-items coords are in the upstream frame (6811-6867 vs local 7028-7084).\par\medskip
\noindent\textit{Allowed standard prerequisites (editorial).}\begin{itemize}
\item Explicit prerequisite boundary. The cited inputs are prerequisites; the argument in the retained span is not replaced by any cited result.
\end{itemize}\medskip
\ifdefined\hypertarget\hypertarget{sourcepointer-1}{}\fi
\noindent\textit{Source reference.} Original lemma 0880: sources/stacks-coherent/pinned-original-source.tex, source lines 6631--6634. Cited by the retained statement or proof but not part of the collected spans; supplied only as a navigation pointer. The argument does not depend on its content.\par\medskip
\ifdefined\hypertarget\hypertarget{sourcepointer-2}{}\fi
\noindent\textit{Source reference.} Original lemma 010I: sources/stacks-coherent/pinned-original-source.tex, source lines 6826--6826. Cited by the retained statement or proof but not part of the collected spans; supplied only as a navigation pointer. The argument does not depend on its content.\par\medskip
\ifdefined\hypertarget\hypertarget{sourcepointer-3}{}\fi
\noindent\textit{Source reference.} Original lemma 0882: sources/stacks-coherent/pinned-original-source.tex, source lines 6655--6666. Cited by the retained statement or proof but not part of the collected spans; supplied only as a navigation pointer. The argument does not depend on its content.\par\medskip
\ifdefined\hypertarget\hypertarget{sourcepointer-4}{}\fi
\noindent\textit{Source reference.} Original lemma 0889: sources/stacks-coherent/pinned-original-source.tex, source lines 6704--6720. Cited by the retained statement or proof but not part of the collected spans; supplied only as a navigation pointer. The argument does not depend on its content.\par\medskip
\begin{lemma}
\label{lemma-torsion-hom-ext}
Let $X$ be a Noetherian scheme and let $\mathcal{I} \subset \mathcal{O}_X$
be a quasi-coherent sheaf of ideals. Any object of
$\textit{Coh}(X, \mathcal{I})$ which is annihilated
by a power of $\mathcal{I}$ is in the essential image of
(\ref{equation-completion-functor}).
Moreover, if $\mathcal{F}$, $\mathcal{G}$ are in $\textit{Coh}(\mathcal{O}_X)$
and either $\mathcal{F}$ or $\mathcal{G}$ is annihilated by a power of
$\mathcal{I}$, then the maps
$$
\xymatrix{
\Hom_X(\mathcal{F}, \mathcal{G}) \ar[d] &
\Ext_X(\mathcal{F}, \mathcal{G}) \ar[d] \\
\Hom_{\textit{Coh}(X, \mathcal{I})}(\mathcal{F}^\wedge, \mathcal{G}^\wedge) &
\Ext_{\textit{Coh}(X, \mathcal{I})}(\mathcal{F}^\wedge, \mathcal{G}^\wedge)
}
$$
are isomorphisms.
\end{lemma}

\begin{proof}
Suppose $(\mathcal{F}_n)$ is an object of $\textit{Coh}(X, \mathcal{I})$
which is annihilated by $\mathcal{I}^c$ for some $c \geq 1$. Then
$\mathcal{F}_n \to \mathcal{F}_c$ is an isomorphism for $n \geq c$.
Hence if we set $\mathcal{F} = \mathcal{F}_c$, then we see that
$\mathcal{F}^\wedge \cong (\mathcal{F}_n)$. This proves the first assertion.

\medskip\noindent
Let $\mathcal{F}$, $\mathcal{G}$ be objects of $\textit{Coh}(\mathcal{O}_X)$
such that either $\mathcal{F}$ or $\mathcal{G}$ is annihilated by
$\mathcal{I}^c$ for some $c \geq 1$. Then
$\mathcal{H} = \SheafHom_{\mathcal{O}_X}(\mathcal{G}, \mathcal{F})$
is a coherent $\mathcal{O}_X$-module annihilated by $\mathcal{I}^c$.
Hence we see that
$$
\Hom_X(\mathcal{G}, \mathcal{F}) =
H^0(X, \mathcal{H}) =
\lim H^0(X, \mathcal{H}/\mathcal{I}^n\mathcal{H}) =
\Mor_{\textit{Coh}(X, \mathcal{I})}
(\mathcal{G}^\wedge, \mathcal{F}^\wedge).
$$
see Lemma \ref{lemma-completion-internal-hom}.
This proves the statement on homomorphisms.

\medskip\noindent
The notation $\Ext$ refers to extensions as defined in
Homology, Section \ref{homology-section-extensions}.
The injectivity of the map on $\Ext$'s follows immediately
from the bijectivity of the map on $\Hom$'s.
For surjectivity, assume $\mathcal{F}$ is annihilated
by a power of $I$. Then part (1) of Lemma \ref{lemma-existence-easy}
shows that given an extension
$$
0 \to \mathcal{G}^\wedge \to (\mathcal{E}_n) \to \mathcal{F}^\wedge \to 0
$$
in $\textit{Coh}(U, I\mathcal{O}_U)$
the morphism $\mathcal{G}^\wedge \to (\mathcal{E}_n)$ is
isomorphic to $\mathcal{G} \to \mathcal{E}^\wedge$
for some $\mathcal{G} \to \mathcal{E}$ in $\textit{Coh}(\mathcal{O}_U)$.
Similarly in the other case.
\end{proof}

\end{document}
